\section{Sistema di Riferimento Affine, Rettifica e Scala Metrica}
\label{sec:affine_rettifica}

Al termine della calibrazione CCalibGA visualizza sul pattern un sistema
di riferimento con assi colorati sovrapposti all'immagine reale. Questa
sezione spiega la geometria sottostante: il pattern fisico porta con sé
un sistema di coordinate affine, la camera lo deforma prospetticamente,
e i parametri stimati permettono di invertire quella deformazione.

\subsection{Sistema di riferimento affine del pattern}
\label{sub:sistema_affine}

Il pattern custom è un oggetto planare fisico: definire su di esso
un'\emph{origine} \(O\) e due assi \(\mathbf{u}, \mathbf{v}\) nel
piano significa dotarlo di un \emph{sistema di coordinate affine}.
Ogni punto del pattern ha coordinate metriche \((X, Y) \in \mathbb{R}^2\)
misure in unità fisiche reali (tipicamente millimetri).

In coordinate omogenee, un punto del pattern si scrive
\(\homo{X}_{\pi} = (X, Y, 1)^T \in \mathbb{P}^2\).
Questo è un punto \emph{proprio}: la terza coordinata è 1, quindi esiste
la sua controparte euclidea finita. Il sistema affine \((O, \mathbf{u},
\mathbf{v})\) non è arbitrario: è imposto dalla geometria del pattern
(dimensione fisica delle celle, posizione dell'origine di rilevamento)
e costituisce il \emph{riferimento metrico} che rende possibile la calibrazione.

\subsection{La camera come proiettività sul piano del pattern}
\label{sub:camera_proiettivita}

Quando la camera inquadra il pattern da una posa generica,
la relazione tra punti del pattern e punti immagine è descritta
dall'omografia \(\mat{H}\):
\[
  \lambda\,\homo{x} = \mat{H}\,\homo{X}_{\pi},
  \qquad
  \mat{H} = \mat{K}\,\bigl[\,\mathbf{r}_1 \;\; \mathbf{r}_2 \;\; \mathbf{t}\,\bigr],
\]
dove \(\mathbf{r}_1, \mathbf{r}_2\) sono le prime due colonne della
rotazione \(\mat{R}\) e \(\mathbf{t}\) è la traslazione della posa.
Poiché \(\mat{H} \in GL(3)\), questa è una proiettività di
\(\mathbb{P}^2\): l'immagine deforma il pattern in modo prospettico,
inclinando e distorcendo le rette parallele del target.

Dal punto di vista della gerarchia delle trasformazioni
(Sezione~\ref{sub:proiettivita}), una proiettività generica
non preserva la struttura affine: rette parallele del pattern
possono convergere nell'immagine verso punti di fuga (punti propri
dell'immagine che corrispondono a punti \emph{impropri} del pattern
portati nel dominio finito dalla proiezione prospettica).

\subsection{La proiettività inversa è ancora una proiettività}
\label{sub:inversa_proiettivita}

Poiché \(\mat{H}\) è invertibile, anche \(\mat{H}^{-1}\) è una matrice
in \(GL(3)\) e quindi descrive una proiettività:
\[
  \homo{X}_{\pi} \sim \mat{H}^{-1}\,\homo{x}.
\]
Questa non è una proprietà banale: significa che per tornare dal piano
immagine al piano del pattern non serve una classe di trasformazione
qualitativamente diversa. L'operazione di \emph{rettifica} dell'immagine
— il warping che riporta il pattern in una vista frontale canonica —
è semplicemente l'applicazione dell'omografia inversa, ancora una
proiettività. OpenCV implementa questo tramite
\texttt{cv2.warpPerspective(img, H\_inv, size)}.

\subsection{Aree nere nella rettifica: punti impropri e regioni non mappate}
\label{sub:aree_nere}

Qualunque utente abbia mai applicato una rettifica prospettica
con OpenCV avrà notato le caratteristiche \emph{zone nere} ai bordi
dell'immagine rettificata. Queste hanno due cause geometriche distinte:

\begin{enumerate}
  \item \textbf{Punti impropri portati nel dominio finito.}
  La proiettività \(\mat{H}\) può mappare punti impropri del piano
  immagine (direzioni, punti di fuga all'infinito) in punti propri del
  piano rettificato, e viceversa: punti dell'immagine originale che
  corrispondono a coordinate del piano rettificato al di fuori
  del rettangolo visibile diventano \emph{non definiti} nel canvas
  di output. Il canvas viene riempito con il colore di background
  (nero per default), generando quelle aree.

  Più formalmente: il canvas di output ha dimensioni finite
  \((W, H)\) in pixel. La rettifica via \(\mat{H}^{-1}\) assegna
  a ogni pixel del canvas una sorgente nel frame originale; se tale
  sorgente cade fuori dal frame originale, o se la terza coordinata
  omogenea del punto risultante è molto piccola (punto vicino
  all'infinito), il pixel non ha colore valido e viene lasciato nero.

  \item \textbf{Distorsione radiale e tangenziale.}
  Prima di applicare l'omografia, OpenCV rimuove la distorsione
  dell'obiettivo tramite i parametri \((k_1, k_2, k_3, p_1, p_2)\).
  Questa undistortion crea localmente aree in cui la griglia di pixel
  si dilata (zone vicino ai bordi per lenti barrel) o si comprime,
  lasciando pixel non coperti che anch'essi appaiono neri.
\end{enumerate}

Le due cause si sommano: in una rettifica completa \emph{undistort +
warpPerspective} le zone nere sono la combinazione di entrambi gli
effetti. Aumentare la dimensione del canvas di output (padding) riduce
le zone nere da prospettiva, ma non quelle da distorsione ai bordi.

\subsection{Correlazione metrica del pattern e importanza della scala}
\label{sub:scala_metrica}

Si chiama \emph{correlazione metrica} del pattern la conoscenza
esplicita delle coordinate fisiche reali dei suoi keypoint.
In CCalibGA ogni keypoint rilevato ha coordinate
\((X_i, Y_i)\) in millimetri nel sistema affine del pattern:
queste informazioni sono usate direttamente dalla funzione
\texttt{cv2.calibrateCamera()} come \texttt{objectPoints}.

Senza la scala metrica, la calibrazione non è determinata:
una camera che osserva un pattern di scala sconosciuta non può
distinguere un pattern piccolo vicino da uno grande lontano
(\emph{ambiguità scala-profondità}). Formalmente, la matrice
di proiezione \(\mat{P} = \mat{K}[\mat{R}|\mathbf{t}]\)
è definita a meno di un fattore di scala globale; fissare la scala
metrica del pattern (es.\ distanza tra keypoint adiacenti = 30 mm)
fissa quella scala e rende la stima di \(\mat{K}\) e di \(\mathbf{t}\)
unica e fisicamente significativa.

\begin{oss}
La scala metrica è importante anche per la visualizzazione del sistema
di riferimento: gli assi proiettati sul pattern (l'asse rosso \(X\),
verde \(Y\), blu \(Z\)) hanno lunghezza fisica nota proprio perché il
sistema affine del pattern ha una scala metrica definita. Se la scala
fosse ignota, gli assi potrebbero essere scalati arbitrariamente e
non avrebbero significato fisico reale.
\end{oss}

In concreto, in CCalibGA la lunghezza fisica delle celle del pattern
custom è parametro di configurazione: cambiarne il valore cambia le
unità di misura di tutta la calibrazione (traslazioni in \(\mathbf{t}\),
focale equivalente se espressa in mm, ecc.). Questo è un vincolo
di calibrazione, non un dettaglio implementativo.

\subsection{Proiezione degli assi: visualizzazione del riferimento affine}
\label{sub:proiezione_assi}

La visualizzazione finale degli assi sul pattern sfrutta
combinazione di tutto quanto descritto:
\begin{enumerate}
  \item Si definisce l'origine \(O = (0,0,0)^T\) e tre punti
  \(P_X = (\ell, 0, 0)^T\), \(P_Y = (0, \ell, 0)^T\),
  \(P_Z = (0, 0, \ell)^T\) nello spazio 3D del pattern
  (con \(\ell\) lunghezza fisica degli assi).
  \item Tramite \texttt{cv2.projectPoints()} — che calcola
  \(\lambda\homo{x} = \mat{K}[\mat{R}|\mathbf{t}]\homo{X}\) —
  questi quattro punti vengono proiettati nel piano immagine.
  \item Si tracciano i segmenti
  \(\overrightarrow{O'P_X'}\), \(\overrightarrow{O'P_Y'}\),
  \(\overrightarrow{O'P_Z'}\) con colori RGB.
\end{enumerate}

Geometricamente, l'asse \(Z\) punta verso la camera (normale al
piano del pattern) e nella proiezione prospettica la sua lunghezza
apparente dipende dall'angolo di inclinazione; è un esempio concreto
di come la proiettività non preservi le lunghezze (a differenza delle
affinità, che non conservano le lunghezze ma conservano il parallelismo).
